\section{贡献一：Intent-Conditioned Controllable Dialogue}

如果直接让语言模型根据完整状态自由续写，它既要决定“应该做什么”，又要决定“应该怎样说”，还可能泄露私有计划或编造非法行动。Cicero 先由战略模块决定近程意图，再让语言模型专注于把意图表达成自然、合时宜的消息。

\subsection{为什么需要显式 intent}

贡献总览 slide 的第一项是 controllable dialogue：模型不仅 condition on game state 和 dialogue history，还 condition on intended actions。这个接口把“策略正确性”和“语言实现”分成可单独评估的两层。

\lecturefigure{slide-25-main-contribution-dialogue.jpg}{Main contribution：用 intended actions 控制对话生成}{官方录像恢复幻灯片 V025；Stanford CS25 Lecture 14}

\begin{knowledgebox}{Intent 在这里不是抽象心理状态}
Cicero 的 intent 主要是可执行的近程行动，例如某个单位本回合或下一回合准备移动、支持或防守到哪里。它比一句自然语言更结构化，便于检查与 board state 一致；但它也比长期战略目标窄，这一限制将在后文返回。
\end{knowledgebox}

\subsection{怎样从人类对话中构造 intent 标签}

人类日志没有手工标注“这句话背后的计划”。系统先训练 intent model，根据 state、history 与对话预测说话者和接收者在回合结束时的 actions；再用实际 orders 为历史消息自动生成近程意图标签。为了减少把谎言当作训练目标，流程会过滤高概率不真实的 turn。

\lecturefigure{slide-26-intent-model-training.jpg}{Intent model training：从对话和最终行动反推消息对应的计划}{官方录像恢复幻灯片 V026；Bakhtin et al., Science 2022}

\lecturefigure{slide-27-intent-annotation.jpg}{Intent annotation：为人类消息自动附加 speaker 与 recipient 的行动意图}{官方录像恢复幻灯片 V027；Stanford CS25 Lecture 14}

\teachervoice{老师特别说明，训练标签来自较可信的 turns，而不是假设所有人类消息都真实。这个选择会牺牲一部分语言多样性，却让 controllability 问题更清楚：先学会稳定表达真实计划，再讨论更复杂的欺骗。}

\subsection{推理时怎样使用 intent}

在线运行时，planner 先给出 Cicero 希望自己和对方采取的行动；intent model 将其组织为条件；dialogue model 再结合 state 和 history 生成多个消息候选。这样，语言模型不必从海量可能策略中重新发明目标，而是解决一个更可控的条件生成问题。

\lecturefigure{slide-28-dialogue-inference.jpg}{Dialogue inference：规划结果先成为 intent，再约束消息候选}{官方录像恢复幻灯片 V028；Stanford CS25 Lecture 14}

若 $m$ 是待发送消息，$z_s$ 和 $z_r$ 分别表示 speaker 与 recipient 的近程行动意图，可以写成
\begin{equation}
p(m\mid s,h,d,z_s,z_r).
\end{equation}
与只建模 $p(m\mid d)$ 相比，这个条件分布显式要求消息同时符合局面、历史与计划。它并不保证生成内容真实或有效，但为一致性检查提供了结构化抓手。

slide 29 的地图示例把 intent 和 message 并排展示。读图时要先看动作箭头，再看语言如何把动作转化为请求或承诺：同一组 orders 可以有不同措辞，但消息若要求一个与箭头冲突的支援，就应被判为 plan-inconsistent。

\lecturefigure{slide-29-intent-conditioned-dialogue.jpg}{Intent-conditioned dialogue：从地图行动意图到自然语言请求}{官方录像恢复幻灯片 V029；Stanford CS25 Lecture 14}

\subsection{结果：可控性与质量为何同时提升}

slide 30 比较三种输入：纯 language model、加入 game-state grounding、再加入 intent grounding。四列分别是与 state 一致、与 plan 一致、人工 high-quality 比例和 perplexity。完整 Cicero 达到 87.30\%、92.86\%、37.30\% 与 7.70，四项均优于前两种设置。

\lecturefigure{slide-30-intent-quality-results.jpg}{Intent grounding 同时改善 state consistency、plan consistency、质量与 perplexity}{官方录像恢复幻灯片 V030；Bakhtin et al., Science 2022}

\begin{knowledgebox}{Perplexity 的第一次解释}
Perplexity 是 token-level cross-entropy 的指数：$\mathrm{PPL}=\exp(H)$。直觉上，它表示模型在每一步面对多少个“等效选择”；越低通常说明对 held-out human dialogue 的概率分配更集中。但 perplexity 不是战略收益，必须与一致性、人工质量和实际游戏结果一起看。
\end{knowledgebox}

\begin{importantbox}{为什么加约束反而让语言更好}
Intent conditioning 缩小了生成问题：模型无需一边发明策略、一边组织语言，只需在给定目标下选择合适表达。因此约束不只是安全护栏，也是一种任务分解，能减少 nonsense 并提高语言模型对有限容量的利用效率。
\end{importantbox}

\subsection{本章小结}

Intent-conditioned dialogue 把战略决策从语言表面实现中分离出来。自动标注提供训练信号，结构化 intent 提供控制接口，而结果表说明这一分解不仅提高 plan consistency，也改善整体语言质量。

